% Clase 7 --- Energía potencial de dos cargas y la convención U(inf) = 0 (§5.3).
% Sólo las DIFERENCIAS de U están definidas; la constante aditiva se fija
% eligiendo U = 0 en el infinito, que es lo que hace la curva pasar por cero.
\begin{center}
\begin{tikzpicture}
\begin{axis}[
  fisfig,
  width=11.5cm, height=6.2cm,
  xlabel={$r$}, ylabel={$U(r)$},
  xmin=0, xmax=4.3, ymin=-1.5, ymax=1.5,
  axis lines=middle,
  xtick=\empty, ytick=\empty,
  every axis y label/.style={at={(ticklabel* cs:1.0)}, anchor=south east},
  every axis x label/.style={at={(ticklabel* cs:1.0)}, anchor=north west},
  legend pos=north east,
]
  % cargas del mismo signo: U > 0, repulsivo
  \addplot[figred, line width=1.2pt, domain=0.42:4.25, samples=140] {0.5/x};
  \addlegendentry{$Q_0Q > 0$ (se repelen)}
  % cargas de signo opuesto: U < 0, ligado
  \addplot[figblue, line width=1.2pt, domain=0.42:4.25, samples=140] {-0.5/x};
  \addlegendentry{$Q_0Q < 0$ (se atraen)}

  % la asíntota: U -> 0 en el infinito
  \draw[figgray, dashed, line width=0.6pt] (axis cs:0,0) -- (axis cs:4.25,0);
  \node[figlblaux, anchor=north east] at (axis cs:4.2,-0.12) {$U \to 0$ cuando $r \to \infty$};
\end{axis}
\end{tikzpicture}\\[0.5em]
{\small\itshape $U = \dfrac{Q_0Q}{4\pi\varepsilon_0 r}$ está definida
\textbf{a menos de una constante aditiva}: sumarle un número a las dos curvas no
cambia ninguna diferencia de energía, que es lo único con sentido físico. La
convención $U=0$ en el infinito es la que las hace tender a cero, y es cómoda
porque separa los dos casos por el signo: $U>0$ significa que hay que
\textbf{gastar} energía para acercarlas, $U<0$ que hay que gastarla para
\textbf{separarlas}.}
\end{center}
